\section{Discussion and Conclusion}
\label{sec:conclusion}

We encoded five principles of distributive justice as an explicit constraint system over
a multi-aircraft noise-routing problem, solved it with an optimising SMT solver, and
measured what each principle does. Five lessons generalise beyond our instance.

\paragraph{The encoding is part of the ethics} Equal admits both an optimisation reading
(``as equal as possible'') and a satisfaction reading (``no two communities differ by
more than $X$''); ordinary usage does not distinguish them, and taking the second --- the
cheaper, more composable choice --- left Equal returning answers interchangeable with
Sufficientarian's wherever the cap does not bind, so what looked like a choice between
principles was a choice between tie-breaks. Two conditions carry that result: no
objective term on either principle, and an optimum lying inside the tighter feasible set.
The first fails under the optimisation reading, which is why the collapse indicts the
reading rather than the principle. The second is empirical, but the underlying inclusion
is not specific to noise --- a bound on the \emph{spread} of any non-negative quantity is
implied by a bound on its \emph{level} --- so any system offering both a comparative and
an absolute fairness cap has one that can only ever be the weaker. A system offering a
menu of fairness principles owes a demonstration that the menu items are distinguishable,
and that is a claim about the instance, not a property of the code.

\paragraph{Ask whether a principle binds, not whether it is implemented} Whether a
principle does anything is a question about the operating point, not about the principle.
At the settings our system ships with, the caps sit at a third of the capacity they
allow and the protected set carries no traffic at all, and three of the five principles
then return the no-fairness routing or something indistinguishable from it. Move the cap
to where it binds and put the protected set where aircraft fly, and the same three
principles separate from that routing by margins a reader can act on.

Measuring which of those regimes one is in is cheap, and it is not the ablation we
started with. Re-solve the same configuration under permutations of the aircraft list:
the problem is unchanged, so the spread over those runs is the solver's indifference, and
a principle whose effect is smaller than that spread is not doing the work being
attributed to it. Running that test is how we found that our own deployed setting was the
wrong one to draw conclusions from.

\paragraph{An unproven answer is not a weak answer, it is a different kind of thing}
Four of our conclusions were wrong, all for the same reason: each compared a timed-out
incumbent against a proven optimum. At the deployed budget Equal and Sufficientarian
appeared to diverge at $N=30$ (arrival 149 against 147); a cap appeared to bind hard at
$N=32$ (158 against 88); Utilitarian appeared to fly more than the others over loud
communities; and the decibel-valued principles appeared to detour substantially from the
shortest path while the others did not. Given four hours, the first pair proves out
identical at 84 --- the earlier answers were $77\%$ from optimal --- the second at 88 with
no binding, the third vanishes, and the fourth reverses: wherever budgets are matched
($N \le 16$) \emph{all five} principles fly every aircraft on a shortest path, and the
apparent detour gap was an artefact of comparing four-hour rows for three principles
against \SI{600}{\second} rows for the other two. Every error pointed in the direction
that made the system look more interesting than it is, none was visible without
re-running, and none would have been caught by inspecting the routings.

We built a standing audit that grades each claim by the status of the rows beneath it,
and it caught the first two. It missed the last two, because it graded by \emph{status}
and the flaw was in \emph{budget}: a table can compare rows that are all `partial' and
still be confounded if some were given six hundred seconds and others four hours. The
audit now checks budget homogeneity within a comparison as well. We report this because
the failure is instructive --- an audit is only as good as the confound it was designed
against, and ours had to be taught a second one.

\paragraph{The metric carries a distributional judgement the principle does not}
Because sound adds in energy, expressing any fairness criterion in \emph{added}
decibels makes already-loud communities roughly $70\times$ cheaper to overfly than quiet
ones at our source level. Every principle we implemented inherits that gradient, and the
one optimising summed decibels most directly --- Utilitarian --- exploits it most, flying
least of all five over the quiet \SI{45}{\dBA} communities ($9.1\%$ of its flight time
against Equal's $34.1\%$) and displacing that traffic onto the mid-level ones. It does
\emph{not} fly more over the loudest tier: that share is pinned at $27.3\%$ for every
principle by where the vertiports sit. The gradient shows up as an exodus from the quiet,
not an influx to the loud. Nothing in the philosophical statement of utilitarianism says
``prefer communities that are already suffering''; the choice of measure says it. Since existing
noise burdens in the United States fall disproportionately on lower-income and minority
neighbourhoods~\cite{casey2017noise,simon2022sociodemographic}, a change-based metric
quietly ratifies the existing distribution. We do not claim an absolute-level metric is
correct; we claim the choice is normative and should be made explicitly.

\paragraph{Tractability is not neutral either} The principles expressible as hard caps
solved to proven optimality in under \SI{200}{\second} at fleet sizes where the two
decibel-valued optimisation principles could not prove optimality at all. If a
deployment picks its fairness rule partly on what it can compute in the time available
--- and a participatory tool with a room full of people waiting certainly does --- then
computational cost is silently doing normative work. Our system's response is to make
the three relevant facts machine-readable and visible: whether a principle was
optimised, whether its optimum was proven, and whether the answer on screen is the
principle's answer or a fallback that ignores it.

\subsection{Limitations and future work}

The acoustic model is deliberately coarse: a tile is treated as uniformly exposed over an
area several kilometres across, and aircraft are point sources with a single level, no
directivity, no altitude and no psychoacoustic weighting~\cite{nasaUAMnoise2020}. Absolute
decibel values should be read as ordinal. Our results concern how fairness rules order
routings under a given noise model, not how loud a real corridor would be.

Proposition~\ref{prop:collapse} holds for any instance, but the empirical findings are
measurements on two instances of one city. Our five principles are also five readings of
five traditions, fixed in advance by the definitions of
Section~\ref{sec:principles}; a maximin
Prioritarian, a threshold-then-Utilitarian Sufficientarian, or an Equal minimising
variance rather than range are all defensible, and comparing such variants is exactly
what this framework exists for. The most direct next step is to close the gap documented
here --- implement Equal as the minimisation its statement calls for, and measure what it
costs and whether it then separates from Sufficientarian. More broadly, the degeneracy we
observe suggests reporting a principle's whole optimal set rather than one arbitrary
member: the honest output of a fairness query may be a family of routings plus a
statement of what the principle does not determine.
